\section{Scheduling: Co-operative and Pre-emptive}

The core scheduler of Apollo-RTOS is a priority based co-operative scheduler. It allows
the highest priority process to run for as long as it needs, until it gives up
execution, when other processes are allowed to continue. Right out of the box, this
simple algorithm allows for several benefits:
\begin{itemize}
    \item It's very simple to implement.
    \item It eleminates a vast majority of cases of race conditions: there can't be a
        race on shared variable between processes, because one process will never
        preempt another.
\end{itemize}
However, it also means that the processes need to be aware of their runtime context,
and need to give up execution periodically to avoid dead-lock. Process can do it in
two ways:
\begin{itemize}
    \item by lowering its priority, or
    \item by going to sleep.
\end{itemize}

With that in mind, the interface for creating and controlling processes is very simple:

\begin{lstlisting}[style=myagc,]
// register process defined by 'proc_def', and pass param as argument
pid_t reg_proc(const proc_def_t *proc_def, void *param);

// increase the priority
void incr_priority(pid_t pid, priority_t priority);

// decrease the priority
void decr_priority(priority_t priority);
\end{lstlisting}

It doesn't take much to define a process:
\begin{lstlisting}[style=myagc,]
struct proc_def_t {
  string      name;
  priority_t  priority;
  size_t      stk_sz;
  runnable_t  func;       // where runnable_t = void (*)(void*);
};
\end{lstlisting}

The runtime representations of a process is stored in a process descriptor table. The
process descriptor holds important information about the running process. The most
important fields of the descriptor are as below:

\begin{lstlisting}[style=myagc]
struct proc_t {
  priority_t  priority;   // current priority
  byte_t      *stk_ptr;   // current stack pointer
  byte_t      *stack;     // start of the stack area
  size_t      stk_sz;     // size of the stack
  void        *param;     // parameter passed to the process function
  // ... and some other stuff
}
\end{lstlisting}

Before we can explain how the scheduling algorithm works, we need to understand the
stack layout in Arm Cortex-M0. 

\begin{emptythm}
    The stack in Arm architecture starts at a higher address, and grows downwards. 

    When execution is pre-empted by an interrupt, \textit{hardware} pushes a few key registers
    onto the stack in a specific order. From lower address to higher, this order is:
    \begin{lstlisting}[style=mytxt]
       new stack pointer -> |                              
                            |  r0                          
                            |  r1                          
       stack        ^       |  r2                          
       grows        |       |  r3                          
       downwards    |       |  r12                         
                            |  lr      Link register       
                            |  pc      Program counter     
       old stack pointer -> |  psr     Status register     
    \end{lstlisting}
    To properly context switch between processes, we need to store some more registers.
    This is done by the software, so we have control over the order. We decided to go
    with the following order:
    \begin{lstlisting}[style=mytxt]
       new stack pointer -> |                              
                            |  r8
                            |  r9
                            |  r10
                            |  r11
                            |  r4
                            |  r5
                            |  r6
    stack pointer after     |  r7
    hardware pushes regs -> |  lr_intr   Return from interrupt address
    \end{lstlisting}
\end{emptythm}

With that out of the way, we are in a position to discuss how scheduling works.

\subsection{Process creation}

The implementation of \lstinline[style=myagc]|reg_proc| can be summarized as follows:

\begin{lstlisting}[
    style=myagc,
    linebackgroundcolor={ \line{4} \line{7} \line{8} },
    title={\file{src/core/sched.cpp}}
]
pid_t reg_proc(const proc_def_t *def, void *param)
{
  auto [name, priority, stk_sz, proc_func] = *def;
  proc_t *proc = procs.alloc();  // allocate a process descriptor
  proc->param = param;  // store the param for future cleanup

  stack_allocator.acquire(proc, stk_sz, proc_func, param, exit);
  incr_priority(pid, priority);

  // ...
  return pid;
}
\end{lstlisting}

The interesting lines here are 7 and 8. In line 7, we ask our local pool allocator to
allocate us a block of memory of size \lstinline[style=myagc]|stk_sz|, and populate it
with the required parameters.

\begin{emptythm}
    For our convenience, we will interpret a newly created process as the same as a
    process that has been just context switched. So the newly created stack should have
    the proper values of the registers in the proper order as discussed above. Which is
    what line 7 does.
\end{emptythm}

If we look at the acquire function:
\begin{lstlisting}[
    style=myagc, 
    linebackgroundcolor={\lines{10}{14}},
    title={\file{src/core/stack\_allocator.h}}
]
void acquire(proc_t *proc, size_t stk_sz, runnable_t func, void *param,
             void (*ret)())
{
  proc->stack = (byte_t *)pool.alloc(stk_sz);  // allocate the stack
  // ...

  auto stk_ptr = (context_t *)(proc->stack + stk_sz - sizeof(context_t));

  // setup initial stack frame
  stk_ptr->psr     = init_psr;        // 0x01000000 to enable Thumb mode
  stk_ptr->pc      = (uint32_t)func;  // entry to the process
  stk_ptr->lr      = (uint32_t)ret;   // exit from the process
  stk_ptr->r0      = (uint32_t)param; // parameter
  stk_ptr->lr_intr = lr_intr_magic;   // return address from interrupts

  proc->stk_ptr = (byte_t *)stk_ptr;
  proc->stk_sz = stk_sz;
}
\end{lstlisting}

On lines 10-14, we populate the ``context'' of the process. We put the entry of the
process to the \texttt{pc} register, the exit sequence from the process to the
\texttt{lr} register, and the parameter of the process to the \texttt{r0} register. We
also set the return address from all interrupt handlers to the top of the stack.


\subsection{Priority modification}

There is a singleton object \lstinline[style=myagc]|cpu| which stores the currently
running process, and current highest priority process:
\begin{lstlisting}[style=myagc, title={\file{src/core/sched.cpp}}]
volatile struct {
  proc_t *hi_proc = nullptr;    // current highest priority process
  proc_t *curr_proc = nullptr;  // currently running process
} cpu;
\end{lstlisting}
This object acts as the \texttt{NEWJOB} register in the AGC. Whenever a new process is
created, it's priority is compared with the priority of the highest priority process.
If the new process has higher priority, it's written to the
\lstinline[style=myagc]|hi_proc| field in \lstinline[style=myagc]|cpu|.

In \lstinline[style=myagc]|reg_proc| line 8, we call
\lstinline[style=myagc]|incr_priority| for the newly created process. This is the
method that does the above checking:
\begin{lstlisting}[
    style=myagc, 
    title={\file{src/core/sched.cpp}},
    linebackgroundcolor={\lines{6}{7}},
]
void incr_priority(pid_t pid, priority_t priority)
{
  proc_t *proc = procs[pid];
  // ... disable irq
  proc->priority = priority; // set the new priority 
  if (cpu.hi_proc == nullptr || cpu.hi_proc->priority < priority)
    cpu.hi_proc = proc;
}
\end{lstlisting}

Now, when a higher priority process is done with its highest priority critical section,
it should ideally lower its priority, to give lower priority processes a chance to run.
It does so using the \lstinline[style=myagc]|decr_priority| method:
\begin{lstlisting}[
    style=myagc, 
    title={\file{src/core/sched.cpp}},
    linebackgroundcolor={\lines{4}{5}},
]
void decr_priority(priority_t priority)
{
  cpu.curr_proc->priority = priority;
  cpu.hi_proc = procs.max_priority();
  change_proc();
}
\end{lstlisting}

First the priority of the current process is decreased. Then the second highest
priority process is found. Then we call \lstinline[style=myagc]|change_proc| method.
This method first checks if the current process is the highest priority process. If
otherwise, it triggers a \lstinline[style=myagc]|pendsv| interrupt, causing the context
switch to happen.
\begin{lstlisting}[
    style=myagc, 
    title={\file{src/core/sched.cpp}},
    linebackgroundcolor={\lines{4}{5}},
]
void change_proc()
{
  // current process is not the highest priority one
  if (cpu.hi_proc != cpu.curr_proc)   
    reschedule();   // triggers pendsv interrupt 
}
\end{lstlisting}

Now, we are ready to see how context switch is done.

\subsection{Context switch}

When \lstinline[style=myagc]|reschedule()| triggers the \lstinline[style=myagc]|pendsv|
interrupt, the hardware pushes the hardware saved registers onto the current process'
stack. Then the \lstinline[style=myagc]|pendsv_handler| is called. From our previous
discussion on the structure of the stack, we now need to push the software saved
registers onto the stack. 

After that, we swap the stack pointer of the current process with the stack pointer of
the incoming process. We then restore the values of the registers based on the values
stored in the stack of the incoming process. 

The storing and restoring of the registers is done in an assembly function, and the
swapping of the stack pointers is doing in a C++ method:

\begin{lstlisting}[style=myarm, title={\file{src/mpx.s}}]
@@@ push the software saved registers onto the stack
    .macro isave
    push {r4-r7, lr}
    mov r4, r8                  @ Copy from high to low
    mov r5, r9
    mov r6, r10
    mov r7, r11
    push {r4-r7}
    mrs r0, msp
    .endm                       @ Return new thread sp

@@@ pop the software saved registers from the stack
    .macro irestore             @ Expect process sp in r0

    msr msp, r0
    pop {r4-r7}
    mov r11, r7 
    mov r10, r6 
    mov r9, r5 
    mov r8, r4
    pop {r4-r7}
    pop {pc}                    @ pop the lr register into the pc
    .endm                       @ resulting in a jump to the new process

@@@ handler for PendSV interupt (context switch)
    .global pendsv_handler
    .thumb_func
pendsv_handler:
    isave                       @ Complete saving of process state
    bl cxt_switch               @ Choose a new process
    irestore                    @ Restore state for that process
\end{lstlisting}

And the definition of \lstinline[style=myagc]|cxt_switch|:

\begin{lstlisting}[
    style=myagc, 
    title={\file{src/core/sched.cpp}},
    linebackgroundcolor={\line{6} \lines{8}{9}},
]
void *cxt_switch(void *stk_ptr)
{
  if (/* cpu.curr_proc is done */)
    /* release resources */
  else 
    cpu.curr_proc->stk_ptr = (byte_t *)stk_ptr;  // store old stk_ptr

  cpu.curr_proc = cpu.hi_proc;  // current proc is the highest priority
  return cpu.hi_proc->stk_ptr;  // restore new stk_ptr 
}
\end{lstlisting}


\subsection{Process suspension}

One way processes can willingly give up execusion is by decreasing its priority.
Another way is for it to suspend its expecution by going to sleep to be woken up at a
later time. 

The mechanism of waking up a process is quite involved, so we will discuss it in a few
steps. There are two ways a process can wake up:
\begin{itemize}
    \item when going to sleep, the process sets up an ``alarm" that wakes it up after a
        certain period of time, or
    \item the process goes to sleep while waiting for an a particular event to occur.
\end{itemize}
The first way will require us to use the waitlist, while the second one will require on
some synchronization mechanism. So before we discuss resuming of a process, we need to
discuss those two topics. 

For now, let's focus on process suspension. 
\begin{emptythm}
    A runnable process is represented by a positive priority. The simplest way to put a
    process to sleep while maintaining its original priority is to just set it's
    priority to be the negative of it's original priority. 
\end{emptythm}

\begin{lstlisting}[
    style=myagc, 
    title={\file{src/core/sched.cpp}},
    linebackgroundcolor={\line{8}},
]
void sleep(time_t period = 0)
{
  if (period > 0)
    /* set an alarm in the waitlist */
  /* else go to bed until some other process wakes me up */

  // negate the priority to represent ASLEEP state
  decr_priority(-cpu.curr_proc->priority);
}
\end{lstlisting}

Another process that wants to wake up a sleeping process only needs to negate the
priority once more:
\begin{lstlisting}[
    style=myagc, 
    title={\file{src/core/sched.cpp}},
    linebackgroundcolor={\line{2}},
]
void wakeup(pid_t pid) { 
  incr_priority(pid, -procs[pid]->priority); 
}
\end{lstlisting}


\subsection{Waitlist: Pre-emptive scheduling}

The waitlist is implemented as a sorted linked list of entries of the format:
\begin{lstlisting}[
    style=myagc,
    title={\file{src/core/waitlist.cpp}},
    linebackgroundcolor={\line{7}},
]
struct waitlist_t {
  string name;
  time_t remaining = 0;
  runnable_t func = nullptr;
  void *param = nullptr;

  waitlist_t *next = nullptr;  // linked list link
};
\end{lstlisting}

\begin{emptythm}
    If a task \(t_0\) is to be done after \(n_0\) ticks, it is placed between tasks
    \(t_1\) and \(t_2\) such that if \(t_1, t_2\) is to be done after \(n_1, n_2\)
    ticks respectively, then \(n_1 \le n_0 \le n_2\). We put, 
    \[
       \texttt{remaining} = n_0 - n_1
    \]
    Hence, the time until \(t_0\) is run is the sum of the \texttt{remaining} fileds in
    all the tasks before \(t_0\).
\end{emptythm}

This is the same algorithm used to store waitlist tasks in the AGC. This has a few
consequences:
\begin{itemize}
    \item It's expensive to keep a sorted linked list, because insertion is \(O(n)\).
        But we put an upper limit on how many waitlist tasks there can be, namely,
        \lstinline[style=myagc]|N_WAITLIST = 16|, so this linear complexity is
        inconsequential.
    \item By having these remaining values, at every tick update, we only need to check
        if the first task in the list has expired. And when we process the first task,
        all we need to do is remove it from the list, and the list invariant will stay
        intact. We don't need to keep a global counter of how much time has passed for
        each task separately.
\end{itemize}

With that said, we have the register and run implementations of the waitlist:
\begin{lstlisting}[
    style=myagc,
    title={\file{src/core/waitlist.cpp}},
    linebackgroundcolor={\line{3} \line{8} \line{16} \line{18}},
]
void reg(string name, time_t interval, runnable_t func, void *param)
{
  waitlist_t *task = alloc();

  waitlist_t *hd = &waitlist_hd;
  time_t prev = 0;
  while (hd->next != nullptr && prev + hd->next->remaining <= interval) {
    prev += hd->next->remaining;  // accumulate remaining times
    hd = hd->next;
  }

  time_t remaining = interval - prev;

  // insert task right after hd
  if (hd->next != nullptr) 
    hd->next->remaining -= remaining;  // invariant for the next task

  *task = {name, remaining, func, param, hd->next};
  hd->next = task;
}
\end{lstlisting}

In the timer interrupt, after every \lstinline[style=myagc]|update_interval| number of
ticks, the waitlist is checked to see if any task's deadline has expired, and if so,
those tasks with expired deadlines are run. 

\begin{lstlisting}[
    style=myagc,
    title={\file{src/core/waitlist.cpp}},
    linebackgroundcolor={\line{5} \line{13}},
]
void run()
{
  auto head = &waitlist_hd;
  if (head->next != nullptr)
    head->next->remaining -= update_interval;

  while (head->next != nullptr && head->next->remaining == 0) {
    auto task = head->next;
    head->next = task->next;

    auto [name, remaining, func, param, next] = *task;
    dealloc(task);
    func(param); // run the task
  }
}
\end{lstlisting}

Here we see the true beauty of this algorithm: we have to do very little in the
interrupt inhibited situation. As long as the tasks are small themselves, this
algorithm does not put any more overhead to the timer interrupt than the necessary bare
minimum.

\begin{lstlisting}[
    style=myagc,
    title={\file{src/drivers/timer.cpp}},
    linebackgroundcolor={\line{8}},
]
void timer1_handler(void)
{
  // advance the timer

  if (/* if update_interval ticks has passed */) {
    // ... disable_irq
    // setup a new stack
    waitlist::run();
  }

  // ... do other stuff 
}
\end{lstlisting}

\subsection{Wakeup alarm}

With the waitlist implemented, setting up an alarm is straightforward:
\begin{lstlisting}[
    style=myagc, 
    title={\file{src/core/sched.cpp}},
    linebackgroundcolor={\lines{5}{8} \line{11}},
]
void sleep(time_t period = 0)
{
  auto proc = cpu.curr_proc;

  auto default_alarm = [](void *ptr) {
    auto proc = (proc_t *)ptr;
    incr_priority(/* pid of proc */, -proc->priority); // wakeup proc
  };

  if (period > 0)   // setup alarm
    waitlist::reg(proc->name, period, default_alarm, (void *)proc);

  // go to sleep
  decr_priority(-proc->priority);
}
\end{lstlisting}

